% sections/03_dataset.tex

\section{Dataset Overview}

The dataset was sourced from Kaggle \cite{moosavi2019dataset} and is backed by
peer-reviewed academic publications, lending it strong reliability.
Data was collected in real-time from multiple traffic APIs, government agencies,
cameras, and road sensors.

\begin{table}[H]
\centering
\begin{tabular}{@{} l l @{}}
  \toprule
  \textbf{Attribute}  & \textbf{Value} \\ \midrule
  Total Records       & 7,728,394 \\
  Features            & 46 original features \\
  Full Time Range     & February 2016 – March 2023 \\
  Analysis Range      & 2016 – 2022 (2023 excluded — partial year) \\
  Coverage            & 49 US states \\
  Licence             & CC BY-NC-SA 4.0 \\
  Severity Scale      & 1 (minor impact) $\rightarrow$ 4 (major disruption) \\
  \bottomrule
\end{tabular}
\caption{Dataset at a glance}
\end{table}

\subsection*{Critical Note on Severity Definition}

\begin{quote}
  \itshape
  Severity in this dataset does \textbf{NOT} measure physical injury or
  fatality. It measures the duration and extent of traffic flow disruption
  caused by the accident. Severity~1 means a short delay; Severity~4 means
  significant, prolonged road blockage.
\end{quote}

This distinction is crucial: high-severity accidents are those that tie up
traffic for longer, correlating with accident complexity, road type, response
time, and infrastructure---not necessarily injury severity.

\subsection*{Reliability Justification}

\begin{itemize}[noitemsep]
  \item Supported by peer-reviewed publications in ACM SIGSPATIAL (2019)
        \cite{moosavi2019risk}.
  \item Multi-source data collection (traffic APIs, law enforcement, road
        sensors) eliminates single-source bias.
  \item Released under CC BY-NC-SA 4.0 with full provenance documentation.
\end{itemize}

\subsection*{Key Features Used}

\begin{table}[H]
\centering\small
\begin{tabular}{@{} l l p{7.5cm} @{}}
  \toprule
  \textbf{Feature}             & \textbf{Type}    & \textbf{Description} \\ \midrule
  Severity                     & Ordinal (1--4)   & Traffic flow impact level \\
  Start\_Time / End\_Time      & Datetime         & Accident start and clearance time \\
  Weather\_Condition           & Categorical      & Raw weather text (grouped into 7 categories) \\
  Sunrise\_Sunset              & Binary           & Day or night at time of accident \\
  Junction / Signal / Crossing & Boolean          & Road infrastructure at accident location \\
  Start\_Lat / Start\_Lng      & Continuous       & GPS coordinates for spatial mapping \\
  State                        & Categorical      & US state of occurrence \\
  Distance(mi)                 & Continuous       & Length of road extent affected \\
  Duration\_Min                & Derived          & End\_Time $-$ Start\_Time in minutes \\
  \bottomrule
\end{tabular}
\caption{Key features used in the analysis}
\end{table}
